\BourbakiExerciseGroup

\begin{enumerate}
\def\labelenumi{\arabic{enumi})}
\tightlist
\item
  Let X be an infinite set and let F be an ultrafilter on X such that the sets of F have empty intersection. For each \(A \in F\) , let \(V_{A}\) denote the subset \(\Delta \cup (A \times A)\) of \(X \times X\) . Show that, as A runs through F, the sets \(V_{A}\) form a fundamental system of entourages for a uniformity \(\mathcal{U}(\mathfrak{F})\) on X, and that the topology induced by this uniformity is discrete.
\end{enumerate}

*2) On the real line R, carrying the additive uniformity, show that as V runs through the filter of entourages of R the sets V(Z) do not form a fundamental system of neighbourhoods of the set Z of rational integers.* (Cf. § 4, no. 3, Corollary to Proposition 4).

*3) Let V be that entourage of the additive uniformity on R which consists of all pairs \((x, y)\) such that either \(|x - y| \leqslant 1\) or \(xy \geqslant 1\) . Show that V is closed in \(R \times R\) , but that if A denotes the (closed) subset of R consisting of all integers \(n \geqslant 2\) , then \(V(A)\) is not closed in \(R\).*

\begin{enumerate}
\def\labelenumi{\arabic{enumi})}
\setcounter{enumi}{3}
\tightlist
\item
  Show that if a uniform space is a Kolmogoroff space (Chapter I, § 1, Exercise 2) then it is Hausdorff.
\end{enumerate}

¶ 5) a) Let X be a uniform space, U its uniformity; for each entourage V of X, let \(\tilde{V}\) be the subset of \(\mathfrak{P}(X) \times \mathfrak{P}(X)\) consisting of all pairs (M, N) of subsets of X such that both \(M \subset V(N)\) and \(N \subset V(M)\) . Show that the sets \(\tilde{V}\) form a fundamental system of entourages of a uniformity \(\tilde{U}\) on \(\mathfrak{P}(X)\) .

\begin{enumerate}
\def\labelenumi{\alph{enumi})}
\setcounter{enumi}{1}
\item
  On the set \(\mathfrak{P}_{0}(\mathbf{X})\) of non-empty subsets of \(\mathbf{X}\) , show that the topology induced by the topology \(\mathfrak{C}(\tilde{\mathcal{U}})\) induced by \(\tilde{\mathcal{U}}\) is finer than the topology \(\mathfrak{C}_{\Phi}\) (Chapter I, § 2, Exercise 7).
\item
  If X has at least two points and if U is Hausdorff, show that the topology induced by \(\mathfrak{G}(\tilde{\mathcal{U}})\) on \(\mathfrak{P}_{0}(\mathbf{X})\) is never coarser than \(G_{\Omega}\) . The topology induced by \(\mathfrak{G}(\tilde{\mathcal{U}})\) on the set \(\mathfrak{F}(\mathbf{X})\) of non-empty closed subsets of X is finer than the topology induced by \(G_{\Omega}\) on this set if and only if, for each closed subset A of X, the sets V(A) form a fundamental system of neighbourhoods of A as V runs through the set of entourages of X.
\end{enumerate}

*d) Show that on the quotient set X/R defined in Chapter I, § 11, Exercise 18, the quotient topology, and the topologies defined by \(\mathfrak{C}_{\Omega}\) , \(\mathfrak{C}_{\Phi}\) and \(\mathfrak{C}(\tilde{\mathcal{U}})\) (where \(\mathcal{U}\) is the usual uniformity on \(\mathbf{R}^2\) ) are all four distinct.*

